\subsection{Test Pilota}
Prima di avviare ufficialmente le sessioni di test con il campione di 15 utenti, è stato condotto un \textit{test pilota} coinvolgendo un singolo partecipante esterno al progetto. Questa prova preliminare si è rivelata estremamente utile per rifinire la procedura sperimentale, chiarire la definizione delle variabili da misurare e capire esattamente come gestire e comportarsi durante lo svolgimento del test.

Nello specifico, il test pilota è servito a far emergere eventuali ambiguità od oscurità nelle istruzioni e a valutare se le richieste dei compiti e degli scenari fossero chiare e di immediata comprensione per l'utente. Inoltre, ha permesso di verificare in modo pratico la durata dell'intero esperimento e delle singole prove, assicurandosi che i compiti non fossero troppo difficili e i tempi troppo stretti o, al contrario, eccessivamente lunghi.

Proprio grazie a questa prova generale, è stato notato che le tempistiche risultavano un po' lunghe e sono state individuate alcune incomprensioni pratiche, legate soprattutto alla formulazione dei due scenari. Sulla base di questi riscontri, le tracce e le richieste sono state ricalibrate e aggiustate, garantendo così la massima fluidità e l'assenza di intoppi durante le sessioni di test definitive.